% ------------------------------------------------------------------------
% Technical note (NOT a paper): what Lightsaber_X changes with respect to
% the reference simulator (old_ref), and how.
%
% Scope: (1) the non-stationary noise environment, (2) reward generation,
% (3) online controller switching. Deliberately EXCLUDED: any description
% of the decision (bandit) algorithms -- here the "policy" is only an
% abstract box that picks the active controller index once per window.
%
% Build: pdflatex code_modifications_note.tex   (no bibliography needed)
% ------------------------------------------------------------------------
\documentclass[11pt,a4paper]{article}

\usepackage[utf8]{inputenc}
\usepackage[T1]{fontenc}
\usepackage[margin=2.6cm]{geometry}
\usepackage{amsmath, amssymb}
\usepackage{booktabs}
\usepackage[colorlinks=true, linkcolor=blue!60!black, urlcolor=blue!60!black]{hyperref}
\usepackage{xcolor}

% Lightweight unit macros (same convention as main.tex; avoids the siunitx dependency).
\newcommand{\SI}[2]{#1\,\text{#2}}
\newcommand{\si}[1]{\text{#1}}
\newcommand{\num}[1]{#1}

\newcommand{\code}[1]{\texttt{#1}}
% let TeX stretch interword space rather than overflow into the margin when a
% long \code{...} identifier cannot be hyphenated
\emergencystretch=2.5em
\newcommand{\fs}{f_{\mathrm{s}}}
\newcommand{\Thold}{T_{\mathrm{hold}}}

\title{\textbf{Lightsaber\_X: simulation-layer modifications\\ with respect to the reference simulator}\\[4pt]
\large Technical note --- environment, reward, and controller handover}
\author{M.~Grillo}
\date{\today}

\begin{document}
\maketitle

\begin{abstract}
\noindent
This note documents, at implementation level, how the \code{Lightsaber\_X} code base
differs from the trusted reference simulator (\code{old\_ref}) in the three areas that
constitute the online controller-selection study: (i)~the construction of a
\emph{non-stationary} noise environment out of three calibrated stationary regimes;
(ii)~the definition and streaming computation of a scalar \emph{reward} that scores
closed-loop performance on consecutive time windows; and (iii)~the \emph{bumpless
handover} machinery that lets the active angular controller be exchanged at run time
without corrupting the loop state. The decision layer that chooses which controller to
play (the bandit) is deliberately \emph{not} described here; throughout, it is treated as
a black box that returns an arm index $a_t \in \{0,1,2\}$ once per decision window.
All physics of the underlying plant is unchanged with respect to the reference; the
equivalence is established in \code{VALIDATION\_vs\_old\_ref.md} and summarised in
Section~\ref{sec:baseline}.
\end{abstract}

\tableofcontents

% ==========================================================================
\section{Baseline and scope}
\label{sec:baseline}

\code{old\_ref} is the original component-graph simulator
(\code{System}/\code{Mirror}/\code{Beam}/\ldots{} classes) of the Sidles--Sigg
angular-control problem for a Fabry--Perot arm cavity. \code{Lightsaber\_X} is built on
its author's Plant-class refactor (\code{Lightsaber.py}: \code{Plant},
\code{Sensors}, \code{SS\_compensation}, \code{Controller}), which we validated against
\code{old\_ref} element by element. The core plant physics is \emph{identical}:

\begin{itemize}
  \item the actuation path (PUM pitch $\to$ test-mass pitch) and the radiation-pressure
        path (test-mass pitch $\to$ pitch) use the same zero--pole--gain models
        ($k=93.52955$ and $k=2.567652$ respectively, identical zeros/poles);
  \item the suspension transfer functions are bit-identical tabulated data;
  \item the intra-cavity power follows the same Fabry--Perot response
        $P_{\mathrm{cav}} = P_{\mathrm{in}}\, t_{\mathrm{ITM}} \, \big/\,
        \lvert 1-\rho_{\mathrm{ITM}} e^{4\pi i\, \delta L/\lambda}\rvert^{2}$,
        with the cavity length fluctuation $\delta L$ high-passed by an elliptic filter
        \code{ellip(2,\,1,\,140,\,$2\pi\cdot50$,\,'high')};
  \item the optical parameters are unchanged: $P=\SI{705}{W}$, $L=\SI{3994.5}{m}$,
        $R_{\mathrm{ITM}}=\SI{1934}{m}$, $R_{\mathrm{ETM}}=\SI{2245}{m}$,
        $t_{\mathrm{ITM}}=0.014$, $\lambda=\SI{1064}{nm}$, $\fs=\SI{256}{Hz}$.
\end{itemize}

Everything documented below sits \emph{on top} of this validated plant: the reference
simulator runs one stationary noise batch with one fixed controller; \code{Lightsaber\_X}
runs a continuous, months-long, regime-switching environment in which the hard-mode
controller may be exchanged at any decision boundary. A Numba just-in-time port of the
per-sample loop (\code{fast\_engine.py}, Section~\ref{sec:fast}) makes such horizons
computationally feasible; it reproduces the Plant-class per-sample sequence exactly and
is validated against it.

Two soft/hard conventions used throughout: the two mirrors' local pitch coordinates
$(\theta_{\mathrm{ITM}},\theta_{\mathrm{ETM}})$ are mapped to the \emph{soft/hard
eigenbasis} by
\begin{equation}
  \label{eq:eig}
  M_{\mathrm{loc}\to\mathrm{eig}} \;=\; \frac{1}{1+r^{2}}
  \begin{pmatrix} 1 & r\\ -r & 1 \end{pmatrix},
  \qquad
  r=\tfrac12\Big(g_1-g_2+\sqrt{(g_1-g_2)^{2}+4}\Big),
  \quad g_i = 1 - L/R_i ,
\end{equation}
and the \emph{hard} mode --- the one destabilised by radiation pressure --- is the
channel that the exchangeable controller acts on and the reward scores. The soft-mode
controller is common to all variants and is never switched.

% ==========================================================================
\section{Non-stationary noise environment}
\label{sec:env}

\subsection{The three stationary regimes}
\label{sec:regimes}

The reference simulator draws one stationary noise realisation from fixed O3-derived
spectra. \code{Lightsaber\_X} replaces this with a \emph{family of three} calibrated
environmental states (``regimes''), each defined by (a)~a seismic input amplitude
spectral density (ASD) file and (b)~a set of scalar multipliers applied to the default
plant flags (\code{bandit\_simulation.py}):

\begin{table}[h]
\centering
\small
\begin{tabular}{lcccc}
\toprule
Regime & ASD file & seismic & input power & sensing\\
\midrule
$R_0$ (nominal)         & \code{ASD\_R0\_nominal.csv}     & 1.0  & 1.00 & 120.0\\
$R_1$ (high microseism) & \code{ASD\_R1\_high\_micro.csv} & 6.0  & 0.95 & 12.0\\
$R_2$ (extreme seismic) & \code{ASD\_R2\_high\_micro2.csv}& 20.0 & 0.90 & 0.25\\
\bottomrule
\end{tabular}
\caption{Regime definitions. The seismic scale multiplies the ITM/ETM ISI injection;
the sensing scale multiplies the white readout-noise ASDs
($n_{\mathrm{hard}}=3\times10^{-14}$, $n_{\mathrm{soft}}=1\times10^{-13}$ in eigenbasis
units$/\sqrt{\mathrm{Hz}}$). $R_2$ pairs \emph{strong} seismic input with a very
\emph{clean} sensor, $R_0$ the opposite --- so no single fixed controller is best in all
regimes, which is the premise of the study.}
\label{tab:regimes}
\end{table}

The per-regime noise itself is synthesised exactly as the reference plant does it
(\code{Plant.create\_tst\_noise\_from\_top}): in the frequency domain, complex white
Gaussian amplitudes are shaped by (root-PSD $\times$ suspension transfer function) and
transformed back,
\begin{equation}
  \label{eq:synth}
  n^{(j)}(t) \;=\; \fs \cdot \operatorname{irfft}\!\Big[\, \eta(f)\;
  s_j \sqrt{S_j(f)}\; T_j(f) \,\Big],
  \qquad
  \eta(f) = \eta_{\mathrm{re}} + i\,\eta_{\mathrm{im}},\;\;
  \eta_{\cdot} \sim \mathcal N\!\big(0,\, \tfrac12\sqrt{T_{\mathrm{blk}}}\big),
\end{equation}
per noise component $j$ (ITM/ETM seismic through the top-stage TFs, OSEM length/pitch
through their TFs), where $s_j$ is the regime scale of Table~\ref{tab:regimes} and
$T_{\mathrm{blk}}$ the synthesis block length. The six shaped components are summed into
the two injected test-mass streams (ITM, ETM). The laser intensity noise is produced the
same way from the O3 relative-power PSD and enters as
$P_{\mathrm{in}}(t) = P\,\big(1 + \fs\!\cdot\!\operatorname{irfft}[\eta\,\sqrt{S_{\mathrm{RIN}}}]\big)$.
The DC bin is zeroed. This reproduces the reference construction \emph{block by block};
each block is an independent stationary realisation of the regime PSDs.

\subsection{Regime mixture: the weight process $W(t)$}
\label{sec:weights}

Non-stationarity is introduced by blending the three regime streams with a
time-dependent, simplex-valued weight vector
$W(t) = \big(W_0(t), W_1(t), W_2(t)\big)$, $\sum_i W_i = 1$:
\begin{equation}
  \label{eq:blend}
  n(t) \;=\; \textstyle\sum_{i} W_i(t)\, n_i(t),
  \qquad
  P_{\mathrm{in}}(t) \;=\; \textstyle\sum_{i} W_i(t)\, P_{\mathrm{in},i}(t),
  \qquad
  s_{\mathrm{sens}}(t) \;=\; \textstyle\sum_{i} W_i(t)\, s_{\mathrm{sens},i},
\end{equation}
i.e.\ the injected suspension noise, the input power, \emph{and} the readout-noise scale
all follow the same mixture. Crucially, the three per-regime realisations $n_i(t)$ are
generated \emph{coherently}: within each synthesis block the random generator is reset to
the same seed before shaping each regime, and the draw order (power first, then the six
test-mass components) is fixed, so the regimes differ \emph{only} through their spectral
shaping and scales, not through independent randomness. A weight crossing
$W: (1,0,0)\to(0,0,1)$ therefore morphs the statistics of one underlying realisation
rather than cross-fading two unrelated noise records (which would transiently
\emph{reduce} the variance by up to a factor 2 at the midpoint).

The weight process itself (\code{RegimeEnv} in \code{bandit\_experiment.py}) is a
compressed-timescale model of a day at a detector site, with two independent layers:

\paragraph{(a) Slow diurnal drift ($R_0 \leftrightarrow R_2$).}
A ``peaky'' diurnal profile
\begin{equation}
  v_\ast(t) \;=\; \tfrac12 + \tfrac12\,\operatorname{sign}\!\big(\cos(2\pi t/P_{\mathrm d})\big)\,
  \big\lvert \cos(2\pi t / P_{\mathrm d}) \big\rvert^{0.33}
\end{equation}
(period $P_{\mathrm d}$; the exponent $0.33$ makes the profile dwell near its extremes
and cross quickly) is \emph{not} followed continuously. Instead, per diurnal cycle a
Poisson number ($\mu = 20$) of \emph{transition events} is drawn, with start times
normally scattered ($\sigma = P_{\mathrm d}/8$) around the two steep points of the cycle
($P_{\mathrm d}/4$ and $3P_{\mathrm d}/4$). At each event the drift value ramps
\emph{linearly} from its current value to the target $v_\ast(t_{\mathrm e})$ over an
exponentially distributed ramp time (mean \SI{40}{s}); between events it holds constant.
The result is a piecewise-linear, stochastically timed day/night oscillation. A hard
threshold then snaps near-extreme values to the pure regimes ($v > 0.75 \mapsto 1$,
$v<0.25\mapsto 0$), so the environment spends most of its time in a \emph{pure} $R_0$ or
$R_2$ state connected by genuine finite-duration ramps --- transitions are continuous,
not instantaneous.

\paragraph{(b) Microseismic spike events ($R_1$).}
Independently, $R_1$ ``storms'' arrive as a Poisson process (rate: \code{spike\_rate}
events per diurnal period; production 6-month run: 1 per day). Each event is a trapezoid
in the $R_1$ weight: attack, hold, and decay durations drawn exponentially with
configurable means (defaults \SI{30}{s}/\SI{150}{s}/\SI{30}{s}; production 6-month run
uses mean hold $=\SI{1500}{s}$).

\paragraph{Composition.} With $v(t)$ the thresholded drift value and $w_1(t)\in[0,1]$
the spike envelope, the final weights are
\begin{equation}
  \label{eq:W}
  W(t) \;=\; \Big( \, v(t)\,\big(1-w_1(t)\big),\;\; w_1(t),\;\; \big(1-v(t)\big)\big(1-w_1(t)\big) \Big),
\end{equation}
i.e.\ a spike overrides the diurnal state and then hands back to it. The process is
seeded, hence fully reproducible; all its parameters (\code{spike\_rate},
\code{spike\_attack/hold/decay}, $\mu$, ramp mean, jitter, exponent, threshold) are part
of the environment configuration.

\subsection{Months-scale realisation: the disk-backed noise cache}
\label{sec:cache}

At $\fs=\SI{256}{Hz}$ a 6-month horizon is $N \approx 4\times10^{9}$ samples; three
regime streams in RAM would exceed \SI{100}{GB}. Because the environment is
\emph{policy-independent} --- every candidate policy must see the identical noise (common
random numbers) --- the blended streams of Eq.~\eqref{eq:blend} are computed \emph{once}
and stored (\code{bandit\_noise\_cache.py}): the two blended test-mass streams and the
blended input power are written as float32 memory-mapped arrays
(\code{sus\_itm.npy}, \code{sus\_etm.npy}, \code{power.npy}, $\sim$\SI{48}{GB} for six
months), together with a manifest holding the per-second weight schedule $W(t)$, the
sensing-scale schedule, the seed, and all environment parameters (which form the cache
key: changing any parameter forces regeneration). Synthesis proceeds in
\SI{4096}{s} blocks; per block, per regime, one call to Eq.~\eqref{eq:synth} with the
block seed \code{seed + 1000\,b}, then the sample-accurate blend with $W$ interpolated
from the per-second schedule. The only artefact of block-wise synthesis is a short
filter transient at each block boundary, negligible on a months-long horizon (the
slowest reward band settles in $\sim\SI{60}{s}$, cf.\ Section~\ref{sec:reward}).

The white \emph{readout} noise is deliberately \emph{not} stored: it is regenerated at
run time, per decision window, from a fixed seed shared by all policies, as Gaussian
samples of standard deviation
\begin{equation}
  \sigma_{\bullet}(t) \;=\; n_{\bullet}\,\sqrt{\fs/2}\; \cdot\; s_{\mathrm{sens}}(t),
  \qquad \bullet \in \{\mathrm{hard}, \mathrm{soft}\},
\end{equation}
added to the two eigen-mode readouts. Since every policy consumes the same number of
draws per window in the same order, the sensing-noise realisation is \emph{identical
across policies} --- the closed-loop trajectories differ only through the controller
choices.

% ==========================================================================
\section{Reward generation}
\label{sec:reward}

The reference simulator has no notion of reward. \code{Lightsaber\_X} scores closed-loop
performance on consecutive windows of $\Thold$ seconds (production: \SI{100}{s}); the
score is designed to (i)~reflect both residual motion \emph{and} control effort across
the band, (ii)~be computable \emph{online} from streaming data without acausal
processing, and (iii)~be comparable across regimes.

\subsection{Multi-band features}

Two signals are scored, both in the eigenbasis of Eq.~\eqref{eq:eig}: the hard-mode
error signal $e_{\mathrm h}(t)$ (the test-mass pitch projected on the hard mode --- the
controller's input before readout noise) and the hard-mode control effort
$u_{\mathrm h}(t)$ (the eigen-projection of the commanded actuation). Each is split into
eight frequency bands by fourth-order Butterworth filters (one low-pass, seven
band-passes):
\begin{center}
\begin{tabular}{ll}
\toprule
band & edges [Hz]\\
\midrule
LP        & $0$ -- $0.3$\\
BP$_1$    & $0.05$ -- $0.3$\\
BP$_2$    & $0.3$ -- $1$\\
BP$_3$    & $1$ -- $3$\\
BP$_4$    & $3$ -- $10$\\
BP$_5$    & $10$ -- $30$\\
BP$_6$    & $30$ -- $80$\\
BP$_7$    & $80$ -- $100$\\
\bottomrule
\end{tabular}
\end{center}
giving 16 features per time block: the RMS of each filtered signal over consecutive
2-second blocks (512 samples),
$F_i(b) = \operatorname{RMS}\big[x_i^{\mathrm{filt}}\big]_{\text{block } b}$.

\subsection{Streaming (stateful) filtering}
\label{sec:stream}

A defect of an early implementation was that the band filters were re-initialised for
every scoring window, so each window began with a spurious filter start-up transient ---
particularly harmful for the sub-\SI{0.3}{Hz} bands whose settling time is
$\mathcal O(\SI{60}{s})$, comparable to the window itself. The current implementation
(\code{reward\_stream.StreamReward}) keeps the internal state $z_i$ of \emph{every} band
filter \emph{across window boundaries}: one \code{StreamReward} object lives for the
whole closed-loop run and its \code{score()} method continues each \code{sosfilt} from
the previous window's final conditions. The reward signal is therefore filtered as one
continuous stream over the entire (months-long) horizon; window boundaries have no
effect whatsoever on the filtering. Equivalently, low-frequency reward content is
attributed to the window in which it occurs, with no per-window warm-up bias.

\subsection{From features to a scalar score}

The raw score of block $b$ is a negatively weighted log-feature sum,
\begin{equation}
  \label{eq:raw}
  Z(b) \;=\; \sum_{i=1}^{16} w_i \,\log\!\big(F_i(b) + \varepsilon\big),
  \qquad
  R_{\mathrm{raw}}(b) \;=\; -\,Z(b), \qquad \varepsilon = 10^{-30},
\end{equation}
so that \emph{smaller} band-RMS (less residual motion, less control effort) gives a
\emph{larger} score, and each band contributes through its \emph{logarithm}, i.e.\
relative (dB-like) changes are weighted rather than absolute power --- a factor-2
improvement in a quiet band counts as much as a factor-2 improvement in a loud one. The
weights $w_i \in [0,1]$ are fixed numbers stored in
\code{bandit\_reward\_weights\_rel2base.npz}; they were characterised offline relative to
the nominal controller $C_0$ (hence \code{rel2base}) so as to emphasise the bands in
which the controller variants actually differ (the file tags its reference as
\code{base\_controller = C0\_nominal}). The same weight file is used by \emph{every}
script in the repository --- calibration, experiment drivers, figures --- so all reported
numbers share one reward definition.

The \emph{window} score is the mean of its per-2s-block raw scores; a \SI{100}{s} window
therefore averages 50 blocks. In the settled regimes the raw score lies in
$\approx[165, 169]$ (Table of Section 5 characterisation runs), dipping to
$\approx 150$ during regime transitions and storms.

\subsection{Normalisation to $[0,1]$ and calibration}
\label{sec:calib}

The decision layer consumes a bounded reward. Two maps are implemented by the
function \code{apply\_reward} in \code{bandit\_calibrate.py}:
\begin{align}
  \text{logistic (default):}\quad
  \rho &= \Big(1 + e^{-(\bar R_{\mathrm{raw}} - c)/s}\Big)^{-1},
  \quad c = 167.5,\; s = 2.5;\\[2pt]
  \text{linear min--max:}\quad
  \rho &= \operatorname{clip}\!\Big( \tfrac{\bar R_{\mathrm{raw}} - \mathrm{lo}}
        {\mathrm{hi} - \mathrm{lo}},\, 0,\, 1 \Big).
\end{align}
The min--max endpoints are \emph{not} hand-set: \code{bandit\_calibrate.py} runs all nine
steady (regime, controller) cells through the actual fast engine
(\SI{1200}{s} each, first \SI{60}{s} discarded), \emph{plus} one fixed-controller pass
over a short \emph{drifting} environment containing regime transitions and $R_1$ spikes,
and sets $[\mathrm{lo},\mathrm{hi}]$ to the pooled min/max widened by a $10\%$ margin.
The drifting pass matters: the steady cells alone bottom out around raw $\approx 163$,
but real transitions dip to raw $\approx 150$; an endpoint set from steady data only
would clip $\sim$10\% of windows to exactly $0$, and a clipped-flat reward during a
recovering transient is indistinguishable from a settled (very bad) regime --- which was
observed to corrupt change-detection logic downstream. With the widened floor a
recovering transient remains a strictly increasing ramp.

The same calibration run also produces (i)~the $3\times3$ table of mean normalised
rewards per (controller, regime) --- used by the Oracle reference policy and for
plotting, and (ii)~$\hat\sigma$, the median within-cell per-window reward standard
deviation, which downstream algorithms use as their noise scale. By construction the
table's per-regime argmax is the designed diagonal ($C_0$ best in $R_0$, $C_1$ in $R_1$,
$C_2$ in $R_2$).

% ==========================================================================
\section{Online controller switching}
\label{sec:switch}

\subsection{The controller variants}

The exchangeable element is the \emph{hard-mode} feedback filter. All three variants
share the reference controller's skeleton --- the 2019-era LIGO ASC zpk (4 zeros
$-0.3436\pm4.11j$, $-0.7854\pm9.392j$; 4 poles $-78.77\pm171.25j$, $-0.062832$,
$-628.32$; $k=5797.86$), an 8th-order boost stage, an integrator zero at
$-2\pi\cdot0.1$, and two elliptic low-pass sections --- and differ in DC gain, low-pass
corner placement, and integrator leak:

\begin{table}[h]
\centering
\begin{tabular}{lcccc}
\toprule
Variant & DC gain & LP1 corner & LP2 corner & integrator leak $f_{\mathrm{leak}}$\\
\midrule
$C_0$ (nominal)      & 30.0  & \SI{8}{Hz}  & \SI{16}{Hz} & \SI{0.03}{Hz}\\
$C_1$ (high micro)   & 44.2  & \SI{10}{Hz} & \SI{20}{Hz} & \SI{0.01}{Hz}\\
$C_2$ (high micro 2) & 50.0  & \SI{12}{Hz} & \SI{26}{Hz} & $0$ (ideal integrator)\\
\bottomrule
\end{tabular}
\caption{Hard-mode controller variants (\code{Controller.set\_feedback\_filter\_hard}).
LP1 = \code{ellip(2,\,1,\,40)}, LP2 = \code{ellip(4,\,1,\,10)} analog prototypes at the
quoted corners; the leak moves the integrator pole to $-2\pi f_{\mathrm{leak}}$. The
analog cascade is discretised by the bilinear transform and factored into second-order
sections (SOS).}
\label{tab:ctrl}
\end{table}

Higher gain and higher corners buy low-frequency suppression (good under strong seismic
input) at the price of more sensing noise re-injected into the loop (bad with a noisy
sensor) --- with the regime table (Table~\ref{tab:regimes}) this makes each variant the
best choice in exactly one regime. The soft-mode filter is identical in all variants and
its state is never touched by a switch.

\subsection{Why switching is delicate}

Each discretised controller is a cascade of $\sim$10 second-order sections holding two
delay states each. Naively ``switching'' by rebuilding the filter (as the reference code
structure would suggest: call \code{set\_feedback\_filter\_hard} with a new name) resets
those states to zero. The consequence is an actuation discontinuity --- the loop, which
was applying an $\mathcal O(\text{settled})$ command, suddenly applies the response of an
empty filter --- i.e.\ an actuator kick whose transient contaminates precisely the
quantity the study measures (the post-switch reward). An early version of this code had
exactly this defect, and the observed ``cost of switching'' was dominated by the
artificial state reset rather than by the controller change itself.

Four handover implementations were therefore built and compared explicitly, in the
script \code{switching\_dynamics\_fast.py} under \code{dev/} (results in the
switching-transient study):
\begin{description}
  \item[cold] the incoming filter starts from zero state (the naive reset);
  \item[hot]  the outgoing filter's state array is transplanted into the incoming filter
              (naive carry-over; the state vectors of different filters are not
              commensurate, so this is also wrong, just differently);
  \item[parallel] the incoming filter has been running \emph{continuously} in the
              background on the same input and is simply made active (its state is
              exactly what it would have been had it always been in charge);
  \item[ramp] parallel, plus a linear cross-fade of the \emph{applied actuation} from the
              outgoing to the incoming filter over a ramp window (bumpless transfer).
\end{description}

\subsection{Production mechanism: parallel banks with a bumpless ramp}
\label{sec:banks}

In the production engine \emph{all three} controller banks are stepped at \emph{every}
sample, each maintaining its own SOS state on the common readout
(\code{fast\_engine.run\_kernel}):
\begin{equation}
  y^{(a)}_k = \mathrm{SOS}^{(a)}\!\big(e_{\mathrm h,k};\, z^{(a)}_k\big),
  \qquad a \in \{0,1,2\},
\end{equation}
so every bank is permanently ``warm''. The applied hard-mode actuation is the active
bank's output; for the first $N_{\mathrm{ramp}}$ samples after a switch
(production: $N_{\mathrm{ramp}} = 2\,\mathrm{s}\cdot\fs = 512$) it is cross-faded from
the previously active bank,
\begin{equation}
  \label{eq:ramp}
  u_{\mathrm h,k} \;=\; (1-\alpha_k)\, y^{(a_{\mathrm{prev}})}_k
                 \;+\; \alpha_k\, y^{(a_{\mathrm{new}})}_k,
  \qquad \alpha_k = \frac{k+1}{N_{\mathrm{ramp}}},
\end{equation}
after which the new bank alone drives the loop. At the switch instant the applied
command is (almost) exactly the old controller's output, and it slews continuously to
the new controller's --- no impulse enters the plant. The slow (Plant-class) engine
implements the same semantics with two filters instead of three
(\code{Controller.switch\_hard}): the outgoing filter is kept running for the ramp
duration and cross-faded out, the incoming filter starts cold but its start-up happens
\emph{inside} the fade, which begins at weight zero.

The measured switching transients under this mechanism are benign: in $R_0$ the
post-switch reward dip is at most $\approx 0.6$ raw-score units even for the largest
gain jump ($C_0 \!\to\! C_2$, gain ratio $1.67$), recovering within a few seconds, and
the integrated loss over \SI{40}{s} is $\lvert \Delta \rvert \lesssim 2$ raw units ---
versus dips an order of magnitude larger for the cold handover. This is what makes
frequent online switching a viable strategy at all, and it is measured, not assumed
(figures \code{fig1\_R0\_transitions\_ramp} and \code{fig6\_handover\_modes} in the
paper).

\subsection{State continuity across the whole horizon}
\label{sec:continuity}

Switching correctness is necessary but not sufficient: the \emph{plant} must also evolve
continuously across decision windows. In the production driver a single persistent state
carries, across every window boundary and hence across the entire multi-month run:
\begin{itemize}
  \item the plant state: previous-sample pitch and applied actuation of both mirrors
        (the loop's one-sample delays), the high-pass filter state of the cavity-length
        path, the radiation-pressure and actuation SOS delay lines;
  \item the Sidles--Sigg compensator state, including the running cavity-power average
        $\bar P$ (updated as $\bar P \leftarrow \bar P + (P_{\mathrm{cav}}-\bar P)/N_c$
        with $N_c$ ramping to $10^{3}$) that sets its gain dynamically;
  \item all three controller-bank states (Section~\ref{sec:banks});
  \item all 16 reward band-filter states (Section~\ref{sec:stream}).
\end{itemize}
A decision window is thus \emph{not} a simulation episode --- it is purely an
accounting boundary at which the reward is read out and the active-arm index may change.
The engine consumes the disk cache in RAM-sized chunks aligned to whole decision windows,
so no window ever straddles a read boundary.

% ==========================================================================
\section{The fast engine}
\label{sec:fast}

None of the above would be usable at months scale with the pure-Python per-sample loop
($\sim 2\times10^{3}$ steps/s). \code{fast\_engine.py} ports the \emph{exact} per-sample
sequence
\begin{center}
  \code{Plant.propagate} $\to$ \code{Sensors.sample\_readout} $\to$
  \code{SS\_compensation.sample\_compensation} $\to$ \code{Controller.sample\_feedback}
\end{center}
--- including both one-sample delays (actuation and Sidles--Sigg compensation are applied
one sample after they are computed) --- into a single Numba-compiled kernel
($\sim 2$--$9\times10^{5}$ steps/s, i.e.\ up to $\sim$3500$\times$ real time). Per sample it:
\begin{enumerate}
  \item forms the beam-spot offsets from the previous pitch, the cavity length change,
        high-passes it, and evaluates the Fabry--Perot power response;
  \item applies the radiation-pressure torque
        $\tau = \tfrac{2}{c}\big(P_{\mathrm{cav}}\,b + d_{\mathrm{off}}(P_{\mathrm{cav}}-P_{\mathrm{dc}})\big)$
        through the rad-to-angle SOS, and the (delayed, negated) commanded actuation
        through the act-to-angle SOS;
  \item adds the injected suspension noise (from the cache) and subtracts the delayed
        Sidles--Sigg compensation, yielding the new pitch;
  \item projects to the eigenbasis and adds the white readout noise
        (Section~\ref{sec:cache});
  \item runs the Sidles--Sigg compensator as a \emph{unit-gain} SOS with a dynamic
        output scale computed from the running power average (exact in steady state);
  \item steps all three controller banks and applies the active/cross-faded actuation
        (Section~\ref{sec:banks}).
\end{enumerate}
All filter coefficients (SOS arrays) and constants are extracted from the validated
Plant-class objects at start-up, so the physics is identical \emph{by construction};
\code{fast\_engine.validate()} reproduces the slow engine's steady-state scores per
(regime, controller) cell. SOS filtering uses the transposed direct-form~II
one-sample step, matching \code{scipy.signal.sosfilt} state conventions exactly, so
states are interchangeable between the two engines.

% ==========================================================================
\section{Summary of differences vs.\ the reference}

\begin{table}[h]
\centering
\small
\begin{tabular}{p{0.24\linewidth}p{0.33\linewidth}p{0.35\linewidth}}
\toprule
 & \code{old\_ref} & \code{Lightsaber\_X}\\
\midrule
Plant physics & component graph & Plant-class refactor; \textbf{validated identical}\\
Noise & one stationary O3 batch & 3 calibrated regimes, coherently synthesised, blended by a stochastic diurnal + Poisson-spike weight process; months-scale disk cache\\
Readout noise & spectral file & white, regime-scaled, regenerated per window from a shared seed (common random numbers across policies)\\
Controller & one fixed zpk & three variants (Table~\ref{tab:ctrl}) as permanently-warm parallel banks; bumpless 2\,s cross-fade handover\\
Reward & none & 16-band log-RMS score of hard-mode error + control effort, streaming stateful filters, calibrated logistic/min--max map to $[0,1]$\\
Loop granularity & one batch = one run & continuous horizon; decision windows are accounting boundaries only (full state continuity)\\
Speed & pure Python & Numba kernel, exact port, $\sim$3500$\times$ real time\\
\bottomrule
\end{tabular}
\caption{One-line summary of the modification set. The decision layer that exploits this
infrastructure is documented separately.}
\end{table}

\end{document}
